\subsection{双接口 \texorpdfstring{$2D$}{2D} 审计律、Toeplitz 非忠实性与判别式坏约化统一}\label{subsec:conclusion-dualinterface-2d-toeplitz-prony-discriminant-badreduction}
沿用定义 \ref{def:pom-hankel-normal-form}、定理 \ref{thm:dephys-hankel-finite-audit}、推论 \ref{cor:xi-hankel-2d-sharpness}、定理 \ref{thm:conclusion-visible-joint-horizon-sharp-2d}、定理 \ref{thm:xi-prony-moment-map-jacobian-delta4}、定理 \ref{thm:xi-toeplitz-metric-spectrum-normal-form}、推论 \ref{cor:xi-negative-spectrum-product-dethankel-square}、定理 \ref{thm:xi-dethankel-vandermonde-square-finite-blaschke}、定理 \ref{thm:xi-hankel-spike-singular-spectrum-separation}、定理 \ref{thm:pom-stiff0-hankel-good-reduction-dim-stability}、定理 \ref{thm:xi-hankel-finitefield-random-completion-nondegenerate} 与定理 \ref{thm:xi-gramshift-bad-primes-weight-discriminant} 的记号。前文已分别在碰撞矩一侧建立了 Hankel 最小实现的 \(2d\) 点锁定，在扫描--Toeplitz 一侧建立了 \(2D\) 视界的 sharp 门槛，并在 Prony Jacobian、Hankel 行列式、负谱乘积与坏素机制之间写出了若干局部闭式。将这些先前分散的接口重新拼接后，可以得到一条此前尚未在结论区单列的刚性链：同一最小 realization 维数 \(D\) 同时支配碰撞接口与扫描接口的最短审计窗；同一 strictification 轨道解释整个 Toeplitz 审计宇宙的非忠实性；同一 Vandermonde--判别式对象既度量微分可逆性，也度量负谱体积与 \(p\)-进坏约化。

\begin{theorem}[双接口的普适 \texorpdfstring{$2D$}{2D} 审计律]\label{thm:conclusion-dualinterface-universal-2d-audit-law}
\leanverified{paper\_conclusion\_dualinterface\_universal\_2d\_audit\_law}
设同一机制同时给出两类有限型接口：
\[
a_n:=S_q(n+2)\qquad(n\ge 0),
\]
其 Hankel 秩为 \(D\)；以及扫描谱指纹
\[
u_n=\sum_{j=1}^{D}c_j r_j^n,
\]
其中 \((r_j,c_j)\) 处于一般位置。则下列三事同时成立：
\begin{enumerate}
  \item 在碰撞接口的最小偏移证书 \(\ell_q\) 处，长度恰为 \(2D\) 的连续窗口
  \[
  (a_{\ell_q},a_{\ell_q+1},\dots,a_{\ell_q+2D-1})
  \]
  唯一锁定全部碰撞动力学，亦即唯一确定 \(\mathrm{HANKELNF}_q\)、递推、特征多项式与 Perron 根。
  \item 在扫描接口上，长度恰为 \(2D\) 的连续前缀
  \[
  (u_0,u_1,\dots,u_{2D-1})
  \]
  唯一锁定全部谱原子与 Toeplitz--PSD 语义。
  \item 长度不超过 \(2D-1\) 的连续样本在两侧都严格不足：碰撞侧仍保留一维 Hankel 纤维，扫描侧仍存在一般位置下的非唯一性见证。
\end{enumerate}
因此，\(2D\) 不是某一单侧接口的偶然门槛，而是同一最小 realization 在“代数最小实现”和“谱最小实现”两侧共同服从的 sharp 信息律。
\end{theorem}

\begin{proof}
碰撞侧由定理 \ref{thm:dephys-hankel-finite-audit} 直接给出：当 Hankel 秩为 \(D\) 时，最小偏移证书 \(\ell_q\) 处的长度 \(2D\) 连续窗口唯一确定 \(\mathrm{HANKELNF}_q\)，因而唯一决定整条碰撞序列及其全部后续指纹。其 sharp 性则由推论 \ref{cor:xi-hankel-2d-sharpness} 给出：任何长度 \(2D-1\) 的连续窗口在一般位置上都仍保留一维自由度，不能唯一恢复全序列。

扫描侧由定理 \ref{thm:conclusion-visible-joint-horizon-sharp-2d} 给出：长度 \(2D\) 的样本前缀既足以完成谱原子唯一恢复，也足以完成 Toeplitz--PSD 的全阶语义判定；而任何长度不超过 \(2D-1\) 的前缀都存在非唯一性见证。于是同一 sharp 阈值在碰撞接口与扫描接口同时成立。证毕。
\end{proof}

\begin{theorem}[Toeplitz--扫描审计函子的非忠实性]\label{thm:conclusion-toeplitz-scan-audit-functor-nonfaithful}
\leanverified{paper\_conclusion\_toeplitz\_scan\_audit\_functor\_nonfaithful}
记 \(\mathscr C\) 为 Carath\'eodory 接口范畴，并对每个对象 \(C\) 赋予核
\[
K_C(w,w'):=\frac{C(w)+C(w')^\ast}{1-ww'^\ast}.
\]
再记 \(\mathscr O\) 为由 \(K_C\) 及其函子性后代生成的审计对象范畴，包括全部 Toeplitz 截断、相应谱铅笔以及由扫描指纹编译出的负向量见证。则自然审计函子
\[
\mathcal F:\mathscr C\longrightarrow \mathscr O
\]
不忠实；更精确地，对任意 \(\eta\in\RR\) 都有
\[
\mathcal F(C+i\eta)=\mathcal F(C).
\]
因此 \(\mathcal F\) 必经 strictification 轨道商 \(\mathscr C/(i\RR)\) 因子化。换言之，若不引入外置寄存轴，则任何仅在 \(\mathscr O\) 内进行的协议都不可能忠实恢复 \(C\) 的全信息。
\end{theorem}

\begin{proof}
定理 \ref{thm:conclusion-toeplitz-gauge-blindness-zero-dimensional-ledger-necessity} 已证明：对任意 \(\eta\in\RR\)，strictification
\[
C\longmapsto C^\sharp:=C+i\eta
\]
不改变核 \(K_C\)，故不改变任何由 \(K_C\) 生成的 Toeplitz 截断、谱读出与 PSD 证书。另一方面，由扫描指纹编译出的负向量见证本身也是从这些 Toeplitz 数据函子性构造出来的，因此同样在 \(C\mapsto C+i\eta\) 下不变。于是 \(\mathscr O\) 中所有可见对象都只能依赖于 \(C\) 的 strictification 轨道，而不能区分同一轨道内的不同代表，故 \(\mathcal F\) 不忠实，并必经商范畴 \(\mathscr C/(i\RR)\) 因子化。证毕。
\end{proof}

\begin{theorem}[Prony 雅可比体积与负谱体积势的精确同一]\label{thm:conclusion-prony-jacobian-negative-spectrum-volume-identity}
\leanverified{paper\_conclusion\_prony\_jacobian\_negative\_spectrum\_volume\_identity}
设扫描谱指纹具有有限原子分解
\[
u_n=4\pi\sum_{j=1}^{\kappa} w_j z_j^n,
\qquad
w_j>0,\quad z_j\in\mathbb D,\quad z_i\neq z_j\ (i\neq j).
\]
令 \(\lambda_1^-,\dots,\lambda_\kappa^-\) 为相应 Toeplitz 审计度量下的负广义特征值，并令
\[
\Phi_\kappa:(\omega_1,\dots,\omega_\kappa,z_1,\dots,z_\kappa)\longmapsto (u_0,\dots,u_{2\kappa-1}),
\qquad
\omega_j:=4\pi w_j,
\]
为 Prony 矩映射。则有严格恒等式
\[
\prod_{i=1}^{\kappa}|\lambda_i^-|
=(4\pi)^\kappa\Bigl(\prod_{j=1}^{\kappa}w_j\Bigr)\cdot \bigl|\det D\Phi_\kappa\bigr|.
\]
亦即，局部可逆性的 Jacobian 体积元与负谱体积势完全同一，仅差一个显式正权因子。
\end{theorem}

\begin{proof}
由定理 \ref{thm:xi-prony-moment-map-jacobian-delta4}，在 \(\omega_j=4\pi w_j\) 的记号下，
\[
\bigl|\det D\Phi_\kappa\bigr|
=(4\pi)^\kappa
\Bigl(\prod_{j=1}^{\kappa}w_j\Bigr)
\Bigl(\prod_{1\le i<j\le \kappa}|z_j-z_i|^4\Bigr).
\]
另一方面，由定理 \ref{thm:xi-dethankel-vandermonde-square-finite-blaschke} 与推论 \ref{cor:xi-negative-spectrum-product-dethankel-square}，
\[
\prod_{i=1}^{\kappa}|\lambda_i^-|
=(4\pi)^{2\kappa}
\Bigl(\prod_{j=1}^{\kappa}w_j^2\Bigr)
\Bigl(\prod_{1\le i<j\le \kappa}|z_j-z_i|^4\Bigr).
\]
将后一式除以前一式即可得到所述恒等式。证毕。
\end{proof}

\begin{corollary}[节点并合的四次判别指数律]\label{cor:conclusion-node-coalescence-quartic-discriminant-law}
沿用定理 \ref{thm:conclusion-prony-jacobian-negative-spectrum-volume-identity} 的设定。设一族参数 \(t\mapsto (w_j(t),z_j(t))\) 光滑变化，并存在一个聚簇 \(I\subset\{1,\dots,\kappa\}\)，满足对所有 \(i,j\in I\),
\[
z_i(t)-z_j(t)=O(\varepsilon(t)),
\qquad
\varepsilon(t)\to 0,
\]
而簇外节点差与全部权重保持远离零。则在一般横截条件下，
\[
\bigl|\det D\Phi_\kappa(t)\bigr|
\asymp
|\varepsilon(t)|^{\,4\binom{|I|}{2}},
\qquad
\prod_{m=1}^{\kappa}|\lambda_m^-(t)|
\asymp
|\varepsilon(t)|^{\,4\binom{|I|}{2}}.
\]
特别地，单对节点并合时，临界指数恰为 \(4\)。
\end{corollary}

\begin{proof}
定理 \ref{thm:conclusion-prony-jacobian-negative-spectrum-volume-identity} 两侧都含有同一 Vandermonde 四次因子
\[
\prod_{1\le i<j\le \kappa}|z_j-z_i|^4.
\]
在簇外差值与权重都保持非退化时，全部会消失的因子都只来自簇内成对差值，因此退化阶精确为
\[
4\cdot \#\{(i,j):i<j,\ i,j\in I\}
=
4\binom{|I|}{2}.
\]
于是两侧都具有同一临界指数。证毕。
\end{proof}

\begin{theorem}[熵缺口趋零时负子空间的秩一塌缩]\label{thm:conclusion-entropy-gap-rankone-negative-collapse}
在 Toeplitz--扫描 Hankel 同构分解
\[
R_N^\ast T_N(C)R_N=
\begin{pmatrix}
P_N & 0\\
0 & -\,\mathsf H_{\kappa-1}^\ast\mathsf H_{\kappa-1}
\end{pmatrix}
\]
的稳定区间内，设
\[
g:=\kappa_{\mathrm{ind}}-S_{\mathrm{def}}\downarrow 0,
\qquad
S_{\mathrm{def}}>0\ \text{固定}.
\]
则 Hankel 负块的奇异谱满足
\[
\sigma_1(\mathsf H_{\kappa-1})=4\pi S_{\mathrm{def}}+O(g),
\qquad
\sigma_j(\mathsf H_{\kappa-1})=O(g)\quad(2\le j\le \kappa).
\]
从而 Toeplitz 审计度量下的负广义特征值满足
\[
\lambda_1^-=
-(4\pi S_{\mathrm{def}})^2+O(g),
\qquad
\lambda_j^-=O(g^2)\quad(2\le j\le \kappa).
\]
故当熵缺口闭合时，全部负方向除主模态外同时二次塌缩；可见缺陷在极限上退化为单一秩一负模。
\end{theorem}

\begin{proof}
定理 \ref{thm:xi-entropy-gap-exponential-suppression-nonzero-fingerprint} 给出
\[
B:=\sup_{n\ge 1}|u_n|\ll g.
\]
又由命题 \ref{prop:xi-defect-entropy-closed-form-invariant} 有 \(u_0=4\pi S_{\mathrm{def}}\)。因此定理 \ref{thm:xi-hankel-spike-singular-spectrum-separation} 立刻给出
\[
\sigma_1(\mathsf H_{\kappa-1})\ge |u_0|-\kappa B=4\pi S_{\mathrm{def}}+O(g),
\qquad
\sigma_j(\mathsf H_{\kappa-1})\le \kappa B=O(g)\ (j\ge 2).
\]
另一方面，
\[
\sigma_1(\mathsf H_{\kappa-1})\le |u_0|+\|\mathsf H_{\kappa-1}-u_0e_0e_0^\top\|_{\mathrm{op}}
=4\pi S_{\mathrm{def}}+O(g),
\]
故首个奇异值确有
\(
\sigma_1(\mathsf H_{\kappa-1})=4\pi S_{\mathrm{def}}+O(g)
\)。

最后，由定理 \ref{thm:xi-toeplitz-metric-spectrum-normal-form}，
\[
\lambda_j^-=-\sigma_j(\mathsf H_{\kappa-1})^2,
\]
代入上述奇异值估计即得所示负广义特征值展开。证毕。
\end{proof}

\begin{corollary}[良素上的随机补全唯一化]\label{cor:conclusion-goodprime-random-completion-unique-extension}
设整数序列 \((a_n)_{n\ge 0}\) 的 Hankel 秩为 \(d\)，并取素数 \(p\) 满足
\[
\mathsf{Stiff}_p=0.
\]
固定任意 \(1\le \delta\le d\)，并仅保留长度 \(2d-\delta\) 的前缀
\[
(\bar a_0,\dots,\bar a_{2d-\delta-1}),
\qquad
\bar a_n:=a_n\bmod p.
\]
将剩余 \(\delta\) 个位置在 \(\FF_p\) 上独立均匀随机补全。则所得长度 \(2d\) 截断以至少
\[
1-\frac{d}{p}
\]
的概率落入 Hankel 非退化开集，因此唯一延拓为一个 Hankel 秩恰为 \(d\) 的全序列。特别地，在一切良素上都存在成功概率 \(1-O(d/p)\) 的 Las Vegas 型有限域最小 realization 重建程序。
\end{corollary}

\begin{proof}
由定理 \ref{thm:pom-stiff0-hankel-good-reduction-dim-stability}，\(\mathsf{Stiff}_p=0\) 蕴含模 \(p\) 约化后的 Hankel 秩仍为 \(d\)。因此真实后缀本身已经给出一个使对应 \(d\times d\) 主块非退化的补全，故补全行列式多项式 \(D(X)\) 非零。于是可直接应用定理 \ref{thm:xi-hankel-finitefield-random-completion-nondegenerate}，得到
\[
\Pr[D(X)=0]\le \frac{d}{p},
\]
也即成功概率至少为 \(1-d/p\)。一旦 \(\det(H_0(X))\neq 0\)，命题 \ref{prop:xi-hankel-window-fiber-delta} 即给出唯一的秩 \(d\) Hankel 延拓。证毕。
\end{proof}

\begin{theorem}[判别式支撑的局部--整体统一律]\label{thm:conclusion-discriminant-support-local-global-unification}
在有限原子指数模型的 \(p\)-整设定下，设特征多项式
\[
\chi_A(T)=\prod_{j=1}^{\kappa}(T-z_j),
\]
权重为 \(w_1,\dots,w_\kappa\)。则同一判别式支撑同时控制：
\[
\text{(a) Archimedean 负谱体积势},
\qquad
\text{(b) non-Archimedean 坏素集合}.
\]
更精确地，
\[
\prod_{i=1}^{\kappa}|\lambda_i^-|
=(4\pi)^{2\kappa}
\Bigl(\prod_{j=1}^{\kappa}w_j^2\Bigr)\cdot
\bigl|\Disc(\chi_A)\bigr|^2,
\]
而在 \(p\)-整模型中，
\[
H_0\bmod \mathfrak p\ \text{可逆}
\quad\Longleftrightarrow\quad
\mathfrak p\nmid
\Bigl(\prod_{j=1}^{\kappa}w_j\Bigr)\Disc(\chi_A).
\]
因此，Archimedean 侧“负谱体积”的消失因子，与 \(p\)-进侧“坏约化”的全部支撑素数，乃由同一判别式对象统一决定。
\end{theorem}

\begin{proof}
由定理 \ref{thm:xi-gramshift-bad-primes-weight-discriminant}，
\[
\det(\mathsf H_{\kappa-1})
=(4\pi)^\kappa
\Bigl(\prod_{j=1}^{\kappa}w_j\Bigr)\Disc(\chi_A).
\]
再由推论 \ref{cor:xi-negative-spectrum-product-dethankel-square}，
\[
\prod_{i=1}^{\kappa}|\lambda_i^-|
=
\bigl|\det(\mathsf H_{\kappa-1})\bigr|^2.
\]
将前式代入即可得到
\[
\prod_{i=1}^{\kappa}|\lambda_i^-|
=(4\pi)^{2\kappa}
\Bigl(\prod_{j=1}^{\kappa}w_j^2\Bigr)\cdot
\bigl|\Disc(\chi_A)\bigr|^2.
\]

另一方面，定理 \ref{thm:xi-gramshift-bad-primes-weight-discriminant} 同时给出：在 \(p\)-整模型下，
\[
H_0\bmod \mathfrak p\ \text{可逆}
\iff
\prod_j w_j\not\equiv 0\pmod{\mathfrak p}
\ \text{且}\
\Disc(\chi_A)\not\equiv 0\pmod{\mathfrak p}.
\]
这正等价于
\[
\mathfrak p\nmid
\Bigl(\prod_{j=1}^{\kappa}w_j\Bigr)\Disc(\chi_A).
\]
故负谱体积与坏素支撑由同一判别式除子统一控制。证毕。
\end{proof}

\endinput
